\subsection{Three additional source families and structural closure}
\label{sec:extension}
We next fixed 30 disjoint combinations of circuits and readouts from YZ-CX,
Hubregtsen, and Kyriienko source modules
\cite{squlearn,hubregtsen2022,kyriienko2021}: ten source circuit
configurations under the same three readouts. The first two modules
use two/four qubits and one/two layers; Kyriienko retains the inventory's
default one-encoding-layer, one-variational-layer construction.
Exact half-integer coefficients from controlled-rotation decompositions
are admitted, without restoring approximate constants to multiples of
$\pi$. Scope, witnesses and resource definitions were fixed before the
run; the nonlinear analysis in Sec.~\ref{sec:nonlinearhub} followed
the observed rank gap.

The bounded Fourier pass completed in 13 settings. Across 252 source
symbol entries it certified 25 exact inactive identities and 195
feasible activity witnesses, leaving 32 entries undecided because
symbolic expansion exceeded its fixed budget. We retain those outcomes
rather than interpreting a budget failure as activity. A separate
structural closure then recovered all 32 remaining inactive entries
using a stronger baseline: global signed Pauli normalization,
common-output isometry deletion, and observable-commutant deletion
that also permits data-dependent rotation angles. All 57 previously
inactive entries in this subset are recovered, and no new inactive
coordinate lies beyond this baseline.

Combining its exact angle-coefficient upper bound with feasible
Jacobian minors closes the rank in 26 settings. The explicit nonlinear
factorization in Sec.~\ref{sec:nonlinearhub} closes one further setting.
At that stage three Kyriienko bounds remained open; activity of every
retained coordinate does not settle these combination ranks. The
post-hoc coefficient-Gram minors and output-gauge proof in
Sec.~\ref{sec:gramkyri} close those three intervals, completing all
30 ranks in this extension and all 54 across the two studies.
For the 26 matching structural bounds, this baseline already attains
the minimum smooth parameter-factor dimension of the tested kernel
family. This does not certify optimal gate, measurement, or shot counts.

\begin{table}[htbp]
\caption{Generic kernel differential rank in the 30-setting extension.
$K$ is the original symbol count. Exact values combine feasible
lower bounds with structural or symbolic upper bounds. Stars mark
the three new coefficient-Gram/gauge closures. The dagger
uses the nonlinear factorization of Eq.~\eqref{eq:hubbeta}.
Rank zero means independence from trainable parameters, not necessarily
a constant function of the inputs.}
\label{tab:extension}
\begin{tabular}{lrrrrrr}
\toprule
Family & $n$ & $L$ & $K$ & $r_{XYZ}$ & $r_Z$ & $r_{\le2}$\\
YZ-CX &2&1&4&3&2&2\\
 &2&2&8&7&6&6\\
 &4&1&8&6&4&6\\
 &4&2&16&14&12&14\\
Hubregtsen &2&1&2&0&2&0\\
 &2&2&4&2&4&2\\
 &4&1&8&8&$2^\dagger$&8\\
 &4&2&16&16&12&16\\
Kyriienko &2&1&6&$4^*$&4&0\\
 &4&1&12&$10^*$&8&$10^*$\\
\bottomrule
\end{tabular}
\end{table}

In all 30 settings, zeroing the additionally Fourier-certified
inactive coordinates after this strengthened baseline produces zero
incremental reduction in Pauli rotation tasks, trainable angle tasks,
or the corresponding two-shift feature-evaluation schedule. These are
ideal operation counts, not executed hardware shots or a fixed-error
shot guarantee. In particular, the nonlinear two-coordinate example
does not constitute another successful internal gate-deletion family.
It instead shows why individual activity and constant-linear null-space
analysis can both miss a nonlinear kernel factorization.
